\section{The resulting infinite family and further specializations}
\label{sec:coverage}

\begin{lemma}[Exponent lifting]\label{lem:lifting}
For $d\geq1$ and $k\geq1$, $\ipra$ proves
\[
 \FLT(d)\longrightarrow\FLT(dk).
\]
\end{lemma}

\begin{proof}
Suppose $x,y,z>0$ and $x^{dk}+y^{dk}=z^{dk}$. The exponent law
$(x^k)^d=x^{dk}$ is proved by quantifier-free induction from the defining
recursion for powers and multiplication. The three positive integers
$x^k,y^k,z^k$ then contradict $\FLT(d)$. Positivity of these powers is
also proved by quantifier-free induction. All variables remain arbitrary,
so the implication is uniform in $d,k$.
\end{proof}

\begin{theorem}[Main exponent family]\label{thm:main}
The theory $\ipra$ proves the single uniform statement
\[
 \boxed{\quad
 \forall n\geq3\quad
 (3\mid n\ \vee\ 4\mid n\ \vee\ 5\mid n\ \vee\ 7\mid n)
 \ \longrightarrow\ \FLT(n).
 \quad}
\]
It also proves the equivalent nontrivial-solution exclusion over signed
integers and over coded rational numbers.
\end{theorem}

\begin{proof}
The four preceding proofs establish the four base exponents. In each
branch of the displayed disjunction, division produces $k=n/d$ with
$n=dk$ and $k\geq1$; apply Lemma~\ref{lem:lifting}.
For even $n$, taking absolute values reduces a signed equation to the
positive case. For odd $n$, write the signed equation as a sum of three
signed $n$th powers equal to zero. The three nonzero terms cannot have
the same sign. Move the term with the exceptional sign to the other side
to obtain a positive equation. These are finite decidable sign cases.
For nonzero rationals, choose the product of their three positive
denominators as a common denominator. Multiplication by its $n$th power
gives a nonzero signed integral solution, which has just been excluded.
All denominator clearing and sign tests are primitive recursive.
\end{proof}

This is an infinite collection of full Fermat specializations, including
all powers of $2$ of exponent at least two, all powers of $3,5,7$, and
every positive multiple of any of the four base exponents. For example,
up to $40$ it proves precisely
\[
 \begin{gathered}
 3,4,5,6,7,8,9,10,12,14,15,16,18,20,21,24,25,\
 27,28,30,32,33,35,36,39,40.
 \end{gathered}
\]
The first exponents $\geq3$ outside this package are
$11,13,17,19,22,23,26,29,31,34,37,38$.
This list describes the reach of the proof, not the truth value of FLT
at those exponents.

\paragraph{Exact coverage.}
The divisibility condition is periodic modulo $420=4\cdot3\cdot5\cdot7$.
The complement has
$3\cdot2\cdot4\cdot6=144$ residue classes: three allowed residues modulo
$4$, two modulo $3$, four modulo $5$, and six modulo $7$. The elementary
Chinese remainder bijection, or a direct scan of $0,\ldots,419$, therefore
gives $420-144=276$ covered classes. Its natural density is
\[
 \frac{276}{420}=\frac{23}{35}\approx65.7142857\%.
\]
There is a completely finite version of this calculation: in every full
block of $420$ consecutive integers exactly $276$ satisfy the divisibility
test. If $N=420q+r$, $0\leq r<420$, the count among $1,\ldots,N$
differs from $(23/35)N$ by less than $420$. Restricting to exponents
$n\geq3$ affects only the initial terms. No analytic limit principle is
needed for this counting estimate.

\begin{corollary}[Mixed exponents]\label{cor:mixed}
Let $a,b,c$ be positive integers. If one of $3,4,5,7$ divides all three
of $a,b,c$, then $\ipra$ proves that
\[
 X^a+Y^b=Z^c
\]
has no positive integral solution. Independently, for all $i,j,k\geq1$
it proves that
\[
 X^{4i}+Y^{4j}=Z^{2k}
\]
has no positive integral solution.
\end{corollary}

\begin{proof}
For the first assertion write $a=d\alpha$, $b=d\beta$, and
$c=d\gamma$ and apply $\FLT(d)$ to
$X^\alpha,Y^\beta,Z^\gamma$. For the second apply
Theorem~\ref{thm:four} to $X^i,Y^j,Z^k$.
\end{proof}

\paragraph{What has not been claimed.}
This paper does not supply a proof of general FLT in $\ipra$. It does
not establish exponent $11$, all regular prime exponents, or an
equivalence between ordinary Lean provability and provability in $\ipra$.
In particular, computing Bernoulli numerators or checking a regularity
condition would not by itself supply an $\ipra$ proof of Kummer's theorem.
We make no claim of independence or unprovability for the omitted cases.
The additional work needed for a new prime exponent is a complete
arithmetic argument with explicit constructions and induction invariants,
comparable to those given here.

